\section{Comparison with an arbitrary aspherical presentation}
\label{sec:presentation}

Let $\Gamma$ be a group with a finite generating alphabet $S$.
Distinct letters are kept distinct, even if they have the same value in
$\Gamma$.  Its generator graph has vertices $g\in\Gamma$ and an oriented
edge $(g,s)$ from $g$ to $gs$ for each $s\in S$.
Put $\Lambda=\Z[\Gamma]$, acting on chains on the left.  For a word $w$
on $S$, write $\bar w$ for its value in $\Gamma$ and $[w]$ for the chain
of the path starting at $1$.  Thus
\[
 \partial_1:\Lambda^{(S)}\longrightarrow\Lambda,
 \qquad \partial_1(e_s)=s-1,
 \qquad \partial_1[w]=\bar w-1.
\]
Here and throughout this section, a parenthesized superscript denotes a
\emph{direct sum}.  All words and all chains have finite support.

\begin{theorem}[The presentation comparison theorem]
\label{thm:presentation}
Suppose that $\Gamma$ is residually finite and that $\mathcal D$ is a
finite indexed list of words on $S$, trivial in $\Gamma$, whose chains
form a free $\Lambda$-basis of $\ker\partial_1$.  Thus the canonical map
\begin{equation}\label{eq:cycle-basis}
 \Lambda^{(\mathcal D)}\longrightarrow\ker\partial_1,
 \qquad e_z\longmapsto[z],
\end{equation}
is an isomorphism.
If $\Gamma$ admits a $K(\Gamma,1)$ of dimension at most two, there is an
$\eta>0$ with the following property.  Assign a label in $H=\SU(2)$ to
every oriented generator edge, assign the inverse label to its reverse,
and assume that every label has distance less than $\eta$ from $1$.
If the product along every translate of every $z\in\mathcal D$ is $1$,
then the product along every closed edge path is $1$.
\end{theorem}

The words in $\mathcal D$ need not present $\Gamma$.  This distinction is
the reason a comparison argument is necessary.  We first arrange an
arbitrary aspherical presentation so that its relators include a finite
list with chains exactly $[z]$.  The remaining relators give independent
word equations in additional loop labels.  A uniform quadratic estimate
isolates the small values of the distinguished relators.  We solve finite
systems on finite group quotients by Lemma~\ref{lem:localized-degree},
and then use compactness to obtain an extension of the original labels.
The presentation used in this argument can have any cardinality.

\subsection{A presentation with distinguished boundary chains}

The correspondence between relator operations and group-ring boundary
operations is classical; compare Fox's transformation~(III) in
\cite[Section~2]{Fox1954}. We give the word realization explicitly
because the presentation here may have arbitrary index sets.

\begin{lemma}\label{lem:aligned-presentation}
Under the hypotheses of Theorem~\ref{thm:presentation}, including the
existence of the two-dimensional $K(\Gamma,1)$, there is a presentation
\[
 \Gamma=\langle S,\mathcal T\mid
        (\rho_z)_{z\in\mathcal D},\ (r)_{r\in\mathcal R}\rangle
\]
with the following properties:
\begin{enumerate}
\item every letter of $\mathcal T$ has value $1$ in $\Gamma$;
\item the complete relator boundary list is a $\Lambda$-basis of the
cycle module of the generator graph on $S\amalg\mathcal T$;
\item $[\rho_z]=[z]$ as edge chains, for each $z\in\mathcal D$;
\item projecting the ordinary boundaries $[r]$, $r\in\mathcal R$, to
the $\mathcal T$-edge coordinates gives an isomorphism
\begin{equation}\label{eq:ordinary-projection}
 \Lambda^{(\mathcal R)}\xrightarrow{\ \cong\ }
 \Lambda^{(\mathcal T)}.
\end{equation}
\end{enumerate}
No finiteness assumption is imposed on $\mathcal T$ or $\mathcal R$.
\end{lemma}

\begin{proof}
Start with any two-dimensional $K(\Gamma,1)$.  Collapse a maximal tree
in its connected one-skeleton and replace the two-cell attaching loops
by finite edge words.  These operations preserve the homotopy type:
a contractible CW subcomplex can be collapsed, and homotopic attaching
maps give homotopy equivalent adjunction spaces
\cite[Propositions~0.17 and~0.18]{Hatcher2002}.  These facts apply to
infinite CW complexes as well; a loop in a graph is homotopic to a finite
edge path.  The resulting presentation complex has one vertex, with
some sets of generators and relators.  Its universal cover is acyclic
and has no cells above dimension two.  Consequently its relator boundary
map is an isomorphism onto the cycle module.  The relator words also
normally generate the kernel of the map from the free generator group
to $\Gamma$.

Adjoin the prescribed letters $S$, with one definition relator for each
letter.  Each new relator boundary has coefficient $1$ on its own new
edge and coefficient $0$ on the other new edges.  Subtracting these
boundaries reduces any cycle uniquely to an old cycle, so the boundary
basis property persists.  For each former generator $a$, choose a word
$w_a(S)$ with value $a$ in $\Gamma$, and replace $a$ by $t_a w_a(S)$.
This is an invertible change of free generators: its inverse sends
$t_a$ to $a w_a(S)^{-1}$ and fixes $S$.  The new letters $t_a$ all have
value $1$.  On edge chains the substitution fixes the $S$-edges and
sends $e_a$ to $e_{t_a}+[w_a]$, so it is a chain isomorphism and preserves
the boundary basis property.  Let $T$ be the set of these new letters,
$R$ the current relator index set, and $I=\mathcal D\amalg T$.
Equation~\eqref{eq:cycle-basis} identifies the cycle module with
$\Lambda^{(I)}$.  In these coordinates the relator boundary map is an
isomorphism
\[
 A:\Lambda^{(R)}\longrightarrow\Lambda^{(I)}.
\]

We next arrange that the prescribed cycles themselves occur as relator
boundaries. An auxiliary copy of the cycle module allows the three
shears below to move this basis into a new relator block.
Adjoin letters $U=(u_i)_{i\in I}$ and defining relators $u_i$.
The stabilized boundary map is
\[
 D_0=\begin{pmatrix}A&0\\0&1\end{pmatrix}:
 \Lambda^{(R)}\oplus\Lambda^{(I)}
 \longrightarrow\Lambda^{(I)}\oplus\Lambda^{(I)},
\]
where the second target summand consists of the $U$-edges.  On its
domain put
\[
 P=\begin{pmatrix}1&A^{-1}\\0&1\end{pmatrix},
 \qquad Q=\begin{pmatrix}1&0\\-A&1\end{pmatrix}.
\]
These are blocks of homomorphisms of left modules; all products below
mean composition.  Direct multiplication gives
\begin{equation}\label{eq:presentation-shears}
 PQP=\begin{pmatrix}0&A^{-1}\\-A&0\end{pmatrix},
 \qquad
 D_0PQP=\begin{pmatrix}0&1\\-A&0\end{pmatrix}.
\end{equation}

We realize each of the successive precompositions $P,Q,P$ by actual
relator words.  A shear leaves one block of words unchanged and adds to
each boundary in the other block a finite sum
$\sum_k n_k g_k[b_k]$ of translates of unchanged boundaries.
Choose a generator word $v_k$ representing $g_k$ and multiply the
changing relator by copies of $v_k b_k^{\pm1}v_k^{-1}$, with the indicated
multiplicities and signs.  All these words represent $1$ in $\Gamma$;
the added boundary is precisely the stated sum.  The unchanged words
belong to both old and new relator lists, so the old changing word is
recovered from the new one by a finite inverse multiplication.  Thus
the old and new lists have exactly the same normal closure.
Every column of $A$ and $A^{-1}$ has finite support because their domains
and targets are direct sums.  Hence this construction uses finite words
for every relator even when the index sets are uncountable.  It does
not require an infinite product of words or a limiting sequence of
presentation moves.

In the final second block, the relator indexed by $i\in I$ has boundary
$(e_i,0)$.  Designate those indexed by $z\in\mathcal D$ as $\rho_z$.
Set $\mathcal T=T\amalg U$.  The remaining second-block relators are
indexed by $T$, and the first-block relators are indexed by $R$; these
are the ordinary relators, with $\mathcal R=T\amalg R$.  Their projected
boundary map is, in these orders,
\begin{equation}\label{eq:ordinary-diagonal}
 \begin{pmatrix}1_T&0\\0&-A\end{pmatrix}:
 \Lambda^{(T)}\oplus\Lambda^{(R)}
 \xrightarrow{\ \cong\ }
 \Lambda^{(T)}\oplus\Lambda^{(I)}.
\end{equation}
Equation~\eqref{eq:presentation-shears} also shows that the complete
boundary list remains a basis.  Normal closure was preserved at every
step, so these are the relators of an actual presentation of $\Gamma$.
\end{proof}

The relators now form an actual presentation, and their boundaries are
aligned with the chosen cycle basis. The next step extracts information
that survives evaluation in the nonabelian group $H$: exact signed
counts will cancel the linear error, leaving a quadratic estimate.

\subsection{Exact counts and a uniform quadratic estimate}

Fix the presentation supplied by Lemma~\ref{lem:aligned-presentation}.
For every $z\in\mathcal D$, choose a finite normal-closure expression
in the free group on $S\amalg\mathcal T$:
\begin{equation}\label{eq:relator-expression}
 z=\prod_{j=1}^{n_z}u_{z,j}\,r_{z,j}^{\sigma_{z,j}}u_{z,j}^{-1},
 \qquad \sigma_{z,j}\in\{1,-1\}.
\end{equation}
Each $r_{z,j}$ is an indexed final relator, ordinary or distinguished.
Let $\mathcal T_0\subset\mathcal T$ consist of the finitely many extra letters
appearing in the words $u_{z,j}$.  For each $t\in\mathcal T_0$, also fix a finite
expression of $t$ as a product of conjugates of final relators and their
inverses.  Finally choose a finite set $\mathcal A\subset\mathcal R$
containing every ordinary relator used in these expressions, both
those for $z$ and those for $t$.

\begin{lemma}\label{lem:exact-counts}
In the free module on all final relator indices,
\begin{equation}\label{eq:translated-counts}
 \sum_{j=1}^{n_z}\sigma_{z,j}\bar u_{z,j}e_{r_{z,j}}
       =e_{\rho_z}.
\end{equation}
\end{lemma}
\begin{proof}
The boundary of a conjugate $u r^{\sigma}u^{-1}$ is
$\sigma\bar u[r]$: the two conjugating paths cancel as chains because
$r$ has group value $1$.  Every factor in
Equation~\eqref{eq:relator-expression} also has value $1$, so multiplying
the factors introduces no further prefix translations in its chain.
Applying the complete relator boundary isomorphism to the left side of
Equation~\eqref{eq:translated-counts} therefore gives $[z]$.
The right side has the same image, namely $[\rho_z]=[z]$.
Injectivity gives the asserted equality, coefficient by coefficient.
\end{proof}

We will use labels either on $\Gamma$ or on any finite quotient $F$ of
$\Gamma$.  In the latter case the graph is formed with the prescribed
images of the generators, retaining distinct edges for distinct alphabet
letters even when their images coincide.  All extra-generator edges are loops in
either graph.  Write $L(w,v)$ for the product of the labels along the
word $w$ starting at $v$, in traversal order.  Endpoints depend only on
the group values of the letters, not on their labels.  Consequently
free word identities can be evaluated on these graphs.  In particular,
\[
 L(u r^{\sigma}u^{-1},v)
   =L(u,v)L(r,v\bar u)^{\sigma}L(u,v)^{-1}
 \quad\text{if }\bar r=1,
\]
with group values projected to $F$ when appropriate.

\begin{lemma}[Uniform estimate]\label{lem:quadratic-estimate}
There is a constant $C\ge1$, depending only on the finite choices above,
with the following property on $\Gamma$ and on every finite quotient.
Suppose that all $S$-labels have distance at most $\epsilon$ from $1$,
and that all translates of the ordinary relators in $\mathcal A$
have product $1$.  Put
\[
 h=\sup_{z\in\mathcal D,\ v}|L(\rho_z,v)-1|,
\]
where the supremum is $0$ if $\mathcal D$ is empty.  If
$0\le h,\epsilon\le1$, then, for every $z\in\mathcal D$ and vertex $v$,
\begin{align}
 |L(z,v)-L(\rho_z,v)|&\le C(h+\epsilon)h,
                      \label{eq:quadratic-estimate}\\
 |L(z,v)-1|&\le C\epsilon.\label{eq:short-word-estimate}
\end{align}
\end{lemma}

\begin{proof}
For unit quaternions, telescoping products and bi-invariance give
\begin{equation}\label{eq:quaternion-estimates}
 \left|\prod_{j=1}^n a_j-\prod_{j=1}^n b_j\right|
       \le\sum_{j=1}^n|a_j-b_j|,
 \qquad
 |ab-ba|\le2|a-1|\,|b-1|.
\end{equation}
The second inequality follows by expanding
$ab-ba=(a-1)(b-1)-(b-1)(a-1)$.
Inverses and conjugation preserve distance from $1$.

Evaluate the fixed expression for each $t\in\mathcal T_0$ at any vertex.
Ordinary relators disappear, and every distinguished factor has
distance at most $h$ from $1$, regardless of its conjugator.  Thus
\[
 |L(t,v)-1|\le M h\qquad(t\in\mathcal T_0)
\]
for one finite constant $M$, independent of the vertex and the group
quotient.  No bound on the conjugators in these auxiliary expressions
is required.  Each $u_{z,j}$ uses only letters of $S\amalg\mathcal T_0$, so another
fixed constant $K$ gives
\[
 |L(u_{z,j},v)-1|\le K(\epsilon+h).
\]
On evaluating Equation~\eqref{eq:relator-expression}, discard its
ordinary factors.  Removing the conjugation from a remaining factor
costs at most
\[
 |aba^{-1}-b|=|ab-ba|
       \le 2K(\epsilon+h)h,
\]
including when the relator occurs with exponent $-1$.
There are at most $n_z$ such factors.

The resulting factors are values of distinguished relators, or their
inverses, at specified vertices.  Any interchange of two adjacent
factors costs at most $2h^2$, by
Equation~\eqref{eq:quaternion-estimates}.  Group together occurrences
having the same relator index and the same vertex; at most
$n_z(n_z-1)/2$ interchanges are needed.  Their signed multiplicities
are prescribed exactly by Equation~\eqref{eq:translated-counts},
translated from $1$ to $v$.  On a finite quotient, coincident vertices
merely add the corresponding coefficients, so the same identity
still applies.  After cancellation within each group, the product is
exactly $L(\rho_z,v)$.  The total error is bounded by
\[
 2n_zK(\epsilon+h)h+n_z(n_z-1)h^2.
\]
Since $\mathcal D$ is finite, this proves
Equation~\eqref{eq:quadratic-estimate} with one $C$.
The word $z$ contains only $S$-letters, so
$|L(z,v)-1|\le\operatorname{length}(z)\epsilon$.
Increasing $C$ proves Equation~\eqref{eq:short-word-estimate} as well.
If the distinguished list is empty, the two assertions are vacuous
and any $C\ge1$ suffices.
\end{proof}

Fix from now on numbers $0<\delta<1$ and $0<\eta<1$ such that
\begin{equation}\label{eq:small-thresholds}
 C(\delta+\eta)<\tfrac12,
 \qquad C\eta<\tfrac{\delta}{2}.
\end{equation}
Under the hypotheses of Lemma~\ref{lem:quadratic-estimate}, with
$\epsilon\le\eta$, the value $h=\delta$ is impossible.  Indeed, taking
the supremum in its two inequalities gives
\begin{equation}\label{eq:no-threshold-crossing}
 \begin{aligned}
 h&\le C\epsilon+C(h+\epsilon)h;\\
 h=\delta&\ \Longrightarrow\
 \delta\le C\eta+C(\delta+\eta)\delta<\delta.
 \end{aligned}
\end{equation}
The supremum need not be attained for this argument.

The constants $C$, $\delta$, and $\eta$ are now fixed, independently of
any finite quotient or finite collection of equations. The exclusion of
$h=\delta$ will keep the degree calculation inside the region of small
distinguished products while the prescribed labels move.

\subsection{Solving finite systems by degree}

\begin{lemma}\label{lem:finite-quotient-solvability}
Let $F$ be any finite quotient of $\Gamma$, and prescribe arbitrary
$S$-edge labels on $F$ of distance less than $\eta$ from $1$.
For every finite set $\mathcal B\subset\mathcal R$ containing
$\mathcal A$, there are labels on all $\mathcal T$-edges such that
\[
 L(r,v)=1\quad(r\in\mathcal B,\ v\in F),
 \qquad
 \max_{z\in\mathcal D,\ v\in F}|L(\rho_z,v)-1|<\delta.
\]
The maximum over an empty distinguished list is interpreted as $0$.
\end{lemma}

\begin{proof}
Tensor the isomorphism~\eqref{eq:ordinary-projection} with the right
$\Lambda$-module $\Z[F]$.  This gives an isomorphism
\begin{equation}\label{eq:quotient-basis}
 \Z[F]^{(\mathcal R)}\xrightarrow{\ \cong\ }
 \Z[F]^{(\mathcal T)}.
\end{equation}
No flatness claim is involved: tensoring the inverse gives an inverse
to the displayed map.  As an abelian-group basis, the source has indices
$(v,r)\in F\times\mathcal R$.  Its image consists of the projected
boundary chains of $r$ based at $v$.  Thus the chains for
$F\times\mathcal B$ are linearly independent over $\Q$.

Give an unknown $x_{v,t}\in H$ to each extra edge $(v,t)$.
In the equation $L(r,v)=1$, the exponent count of each unknown is
exactly its coefficient in the projected boundary chain.  This holds
even though the fixed $S$-labels need not commute: the path traversed
is determined by the generator values in $F$, and all $t\in\mathcal T$
have value $1$.  Write the $m=|F|\,|\mathcal B|$ exponent vectors as
rows.  Their union of supports is finite, and their row rank is $m$.
Choose $m$ unknowns giving a nonsingular $m\times m$ integer minor
$B$, and set every other extra label equal to $1$.

Move all the prescribed $S$-labels from $1$ to their final values
along the short quaternion arcs, with common parameter $\tau\in[0,1]$.
These arcs stay at distance less than $\eta$ from $1$.
The selected equations now define a continuous family of smooth maps
\[
 f_\tau:H^m\longrightarrow H^m.
\]
For an assignment $x\in H^m$, let
\[
 h(\tau,x)=\max_{z\in\mathcal D,\ v\in F}
       |L_\tau(\rho_z,v;x)-1|,
 \qquad \mathcal O_\tau=\{x:h(\tau,x)<\delta\}.
\]
Only finitely many words occur here, so $h$ is continuous.
Every root $f_\tau(x)=\one$ satisfies all equations from $\mathcal A$.
Equation~\eqref{eq:no-threshold-crossing} therefore implies
\begin{equation}\label{eq:root-isolation}
 f_\tau(x)=\one\quad\Longrightarrow\quad h(\tau,x)\ne\delta.
\end{equation}
If $m=0$, the domain and target are points.  The unique assignment
starts with $h=0$, and continuity together with
Equation~\eqref{eq:root-isolation} keeps $h<\delta$ throughout.
This proves the assertion in that case; henceforth assume $m\ge1$.

We justify degree invariance with these moving open sets.  For fixed
$\tau_0$, the roots with $h(\tau_0,x)<\delta$ form a compact set:
by~\eqref{eq:root-isolation} they are also the roots with
$h(\tau_0,x)\le\delta$.  Choose an open neighborhood $V$ containing
exactly these roots, with
$\overline V\subset\mathcal O_{\tau_0}$.
If there are no such roots, take $V=\varnothing$.
Compactness of $H^m$ and continuity show that, for all $\tau$ sufficiently
close to $\tau_0$, every root with $h(\tau,x)<\delta$ lies in $V$,
every root in $\overline V$ has $h(\tau,x)<\delta$, and there is no
root on $\partial V$.  For example, a sequence of selected roots
outside $V$ with parameters tending to $\tau_0$ would have a limit
root with $h\le\delta$, hence $h<\delta$, outside $V$, a contradiction.
The other assertions follow in the same way, or from
$\overline V\subset\mathcal O_{\tau_0}$.
Excision and homotopy invariance of ordinary degree now give
\[
 \deg(f_\tau,\mathcal O_\tau,\one)
       =\deg(f_\tau,V,\one)
       =\deg(f_{\tau_0},V,\one)
\]
near $\tau_0$.  Thus the selected degree is locally constant, and hence
constant on $[0,1]$.

At $\tau=0$, the map $f_0$ is a pure word map with exponent matrix $B$:
all fixed coefficients are $1$.  The set $\mathcal O_0$ is invariant
under simultaneous conjugation, contains the identity assignment,
and its boundary contains no root.  Lemma~\ref{lem:localized-degree}
therefore gives
\[
 \deg(f_0,\mathcal O_0,\one)\ne0.
\]
The same holds at $\tau=1$, so a root exists in $\mathcal O_1$, as
required.
\end{proof}

Finite quotients now provide solutions for every finite set of ordinary
relators, with the distinguished products uniformly bounded. To return
to the original graph, we must preserve its prescribed $S$-labels exactly.
The following compactness argument does this one finite collection of
constraints at a time.

\subsection{Finite-coordinate compactness and the conclusion}

We use compactness of arbitrary products of compact spaces
\cite[Theorem~37.3]{Munkres2000}. Its unrestricted index set is essential:
the set of extra generators in the hypothetical presentation need not
be countable.

\begin{proof}[Proof of Theorem~\ref{thm:presentation}]
Let $a$ be any prescribed labeling of the $S$-edges on $\Gamma$ satisfying
the theorem's hypotheses, with the fixed $\eta$ of
Equation~\eqref{eq:small-thresholds}.  In the compact product
\[
 \mathscr P=H^{\Gamma\times\mathcal T}
\]
consider all the following conditions on the unknown extra-edge labels:
\begin{align}
 L(r,g)&=1 &&(r\in\mathcal R,\ g\in\Gamma),
                           \label{eq:compact-ordinary}\\
 |L(\rho_z,g)-1|&\le\delta
             &&(z\in\mathcal D,\ g\in\Gamma).
                           \label{eq:compact-distinguished}
\end{align}
Each condition is closed and depends on finitely many coordinates.
We verify the finite intersection property, keeping the original
labeling $a$ fixed throughout.

First choose an arbitrary finite collection $\mathscr C$ of the
conditions~\eqref{eq:compact-ordinary}--\eqref{eq:compact-distinguished}.
Let $P\subset\Gamma$ be the finite set of starting vertices of all
\emph{positively oriented} $S$-edges used when evaluating the words
in $\mathscr C$.  An inverse occurrence traversing the edge from
$g$ to $gs^{-1}$ uses the positively oriented edge starting at
$gs^{-1}$, and this starting vertex is included in $P$.
Let $\mathcal B$ contain $\mathcal A$ and all ordinary relator indices
appearing in $\mathscr C$.  These choices are made before choosing a
finite quotient.

Residual finiteness supplies a finite quotient
$q:\Gamma\twoheadrightarrow F$ that is injective on $P$:
separate the finitely many nonidentity elements $p^{-1}p'$ for distinct
$p,p'\in P$, and let $F$ be the image of $\Gamma$ in the product
of the resulting finite quotients.
For each $S$-edge actually used in $\mathscr C$, assign to its image
edge in $F$ its prescribed label from $a$.  Injectivity on $P$ ensures
that no two inconsistent prescriptions are placed on the same edge.
Give every remaining $S$-edge of $F$ label $1$.  All these labels have
distance less than $\eta$ from $1$.

Apply Lemma~\ref{lem:finite-quotient-solvability} to this finite quotient,
this assignment, and $\mathcal B$.  Pull the resulting extra-edge labels
back by $q$, and combine them with the \emph{original} $S$-labels $a$
on $\Gamma$.  In every word from $\mathscr C$, each traversed $S$-edge
has the same label as its image in $F$, by construction; the same is
true for each extra edge by pullback.  Its path product therefore
agrees with the quotient path product.  The selected ordinary equations
hold, and the selected distinguished products even have distance
strictly less than $\delta$.  This supplies a point of $\mathscr P$
satisfying $\mathscr C$.

We have proved: for each finite collection of constraints, one can
choose a finite quotient, copy the finitely many required prescribed
labels, solve a finite system there, and pull back a solution to that
collection.  The quotient may depend on the collection; the original
assignment $a$ is never assumed to descend to any one finite quotient.
Compactness of $\mathscr P$ now gives an extension satisfying
all of~\eqref{eq:compact-ordinary}--\eqref{eq:compact-distinguished}
simultaneously.

For this extension, $h\le\delta$ and all the equations in $\mathcal A$
hold.  The original hypothesis says $L(z,g)=1$ for every $z$ and $g$.
Taking suprema in Equation~\eqref{eq:quadratic-estimate}, with
$\epsilon=\eta$, gives
\[
 h\le C(h+\eta)h\le C(\delta+\eta)h<\tfrac12h
 \quad\text{if }h>0.
\]
Thus $h=0$.  All distinguished relators, as well as all ordinary
relators, have product $1$ at every vertex.
Finally, a closed edge path spells a word representing $1$ in $\Gamma$.
Because Lemma~\ref{lem:aligned-presentation} gives an actual presentation,
that word is a finite product of conjugates of final relators and their
inverses.  Evaluating at the starting vertex makes every factor $1$.
Hence the original closed path has product $1$, as claimed.
\end{proof}
